% Chapter 2

\chapter{安全对齐机制与漏洞机理分析}

\section{大语言模型及安全对齐技术基础}

\section{大语言模型概述}

大语言模型（Large Language Models, LLMs）是一类基于深度神经网络构建的生成式语言模型，通常依托大规模语料进行预训练，以学习自然语言中的统计分布规律与语义表示。相较于传统面向特定任务设计的自然语言处理模型，大语言模型普遍采用“预训练—对齐”的范式：首先通过开放域语料进行无监督或自监督训练，获得通用语言能力；随后通过指令微调、偏好优化等方法对模型行为进行约束，使其适配具体应用场景\cite{touvron2023llama,rafailov2023dpo}。

从模型结构上看，当前主流大语言模型多基于Transformer架构，通过自注意力机制刻画序列中不同位置之间的依赖关系，从而实现对长距离上下文的有效建模。在自回归生成框架下，模型通过最大化条件概率分布来生成文本，即在给定上下文$x_{1:n}$的条件下预测下一个词元$x_{n+1}$的概率\cite{touvron2023llama}：
\begin{equation}
P(x_{n+1} \mid x_{1:n})
\end{equation}
通过对上述条件概率的逐步建模与采样，模型能够生成语义连贯的文本序列。该生成机制决定了模型输出本质上是对训练数据分布的拟合结果，而非基于显式规则的推理过程。

\begin{figure}[htbp]
\centering
\resizebox{0.98\textwidth}{!}{\begin{tikzpicture}[
    x=1pt,
    y=1pt,
    font=\small,
    node distance=16pt,
    >=stealth,
    stage/.style={
        draw=#1!75!black,
        fill=#1!10,
        rounded corners=10pt,
        line width=0.9pt,
        minimum width=106pt,
        minimum height=64pt,
        align=center
    },
    subtag/.style={
        draw=#1!70!black,
        fill=white,
        rounded corners=6pt,
        line width=0.8pt,
        inner xsep=7pt,
        inner ysep=4pt,
        font=\scriptsize
    },
    flow/.style={
        ->,
        line width=1pt,
        draw=gray!75
    },
    contextline/.style={
        ->,
        line width=0.9pt,
        draw=teal!70!black,
        dashed
    }
]
\def\W{730}
\def\H{248}

\fill[gray!4, rounded corners=16pt] (0,0) rectangle (\W,\H);
\draw[gray!35, rounded corners=16pt, line width=1pt] (0,0) rectangle (\W,\H);

\fill[cyan!18, rounded corners=14pt] (18,200) rectangle (712,230);
\node[anchor=west, font=\bfseries\small, text=black] at (34,215)
{Large Language Model Generation Pipeline};
\node[anchor=east, font=\scriptsize, text=gray!70!black] at (694,215)
{context-conditioned autoregressive decoding};

\node[stage=cyan] (prompt) at (78,118) {
    \textbf{输入提示}\\[4pt]
    用户问题\\
    系统指令\\
    历史消息
};

\node[stage=blue] (token) at (198,118) {
    \textbf{分词编码}\\[4pt]
    Text $\rightarrow$ Tokens\\
    位置编码\\
    向量表示
};

\node[stage=teal, minimum width=158pt, minimum height=84pt] (transformer) at (372,118) {
    \textbf{Transformer 堆叠}\\[4pt]
    多层自注意力\\
    前馈网络\\
    残差与归一化
};

\node[subtag=teal] (context) at (372,66) {上下文窗口};
\node[subtag=teal] (attention) at (372,169) {注意力计算};

\node[stage=orange] (prob) at (544,118) {
    \textbf{下一词概率分布}\\[4pt]
    logits\\
    Softmax\\
    候选词概率排序
};

\node[subtag=orange] (sampling) at (544,66) {概率采样};

\node[stage=red] (decode) at (650,118) {
    \textbf{自回归采样}\\[4pt]
    Top-k / Top-p\\
    温度控制\\
    生成下一个词元
};

\node[stage=violet] (output) at (650,118-92) {
    \textbf{输出回复}\\[4pt]
    逐词展开\\
    语义连贯\\
    最终响应
};

\draw[flow] (prompt.east) -- (token.west);
\draw[flow] (token.east) -- (transformer.west);
\draw[flow] (transformer.east) -- (prob.west);
\draw[flow] (prob.east) -- (decode.west);
\draw[flow] (decode.south) -- (output.north);

\draw[contextline] (prompt.north) .. controls (116,182) and (314,188) .. (transformer.north west);
\draw[contextline] (token.north) .. controls (228,180) and (326,180) .. (attention.west);
\draw[contextline] (context.east) .. controls (430,66) and (492,66) .. (sampling.west);

\draw[flow, rounded corners=10pt] (decode.west) |- (592,30) -| node[pos=0.55, below=8pt, font=\scriptsize, text=gray!70!black] {将新词元拼接回上下文并继续预测} (token.south);

\node[font=\scriptsize, align=left, text=gray!75!black] at (108,30) {
    核心机制: 在给定上下文 $x_{1:n}$ 条件下，模型持续预测 $P(x_{n+1}\mid x_{1:n})$
};
\end{tikzpicture}
}
\caption{大语言模型生成机制示意图}
\label{fig:llm-generation-pipeline}
\end{figure}

\begin{equation}
P(x_{1:T})=\prod_{t=1}^{T}P(x_t \mid x_{<t})
\end{equation}
其中，$x_t$表示第$t$个词元，$x_{<t}$表示第$t$步之前的前缀上下文，$T$表示输出序列长度。该式与前文的单步条件概率共同构成了自回归语言模型的完整生成定义。

随着模型参数规模、训练数据规模以及优化方法的不断提升，大语言模型在多任务场景中表现出显著的泛化能力，并在指令理解、复杂推理及多轮对话等任务中取得了优异性能。同时，模型在一定条件下还表现出上下文学习与链式推理等能力，使其能够在无需额外训练的情况下完成多种复杂任务\cite{touvron2023llama}。

然而，大语言模型的能力提升并未伴随安全性的同步增强。预训练阶段的优化目标主要关注语言建模能力，而非对有害内容的识别与约束，因此模型不可避免地继承训练数据中的偏差与不确定性。在缺乏额外约束的情况下，模型可能在特定输入条件下生成不符合安全规范的内容。

此外，大语言模型的生成过程高度依赖输入上下文，其输出行为不仅由当前提示决定，还受到历史上下文的持续影响。这种上下文依赖特性在提升模型表达能力的同时，也为对抗性输入提供了作用空间，使得攻击者可以通过构造特定语境诱导模型偏离既定行为约束。

综上，大语言模型作为一种基于概率建模的生成系统，其行为本质上受数据分布与上下文条件共同驱动。在缺乏有效约束机制的情况下，其输出存在被诱导产生不安全内容的风险。因此，有必要在此基础上进一步引入安全对齐机制，以对模型行为进行系统性约束。

\section{安全对齐技术体系}

在大语言模型的应用过程中，仅依赖预训练阶段获得的语言建模能力难以满足实际场景中的安全要求，因此需要引入额外的约束机制对模型行为进行规范。所谓安全对齐（Safety Alignment），是指通过训练、控制与评估等多种手段，使模型生成内容尽可能符合人类价值、平台规范及应用安全边界。从实现路径上看，安全对齐并非单一方法，而是由模型内部训练机制与外部防护机制共同构成的综合技术体系。

在模型内部层面，安全对齐主要通过后训练（post-training）过程实现，其核心目标是在保持模型能力的同时，引导其生成行为向“安全且有用”的方向收敛。其中，有监督微调（Supervised Fine-Tuning, SFT）构成了最基础的对齐路径。SFT通过构建高质量的指令—响应数据，对预训练模型进行进一步训练，使其在面对潜在风险请求时倾向于输出拒答、警示或安全替代内容。该方法能够有效塑造模型的交互风格与基本行为边界，但其约束能力依赖于训练样本的覆盖范围，对于未见或变形输入的泛化能力相对有限\cite{markov2022holistic,phute2024llmselfdefense}。

在此基础上，基于人类反馈的强化学习（Reinforcement Learning from Human Feedback, RLHF）进一步引入偏好建模机制。该方法通过构建人类偏好数据并训练奖励模型，对模型输出进行质量评估，并利用强化学习算法优化生成策略，使模型在多种候选响应中更倾向于选择符合安全与偏好约束的结果。相较于SFT，RLHF能够在“有用性”与“安全性”之间实现更细粒度的权衡。然而，由于奖励模型本质上仍是对有限偏好数据的近似，其泛化能力受限，且在复杂场景下可能出现奖励偏移或策略失真等问题\cite{rafailov2023dpo,ethayarajh2024kto}。

为降低训练复杂度并提升优化稳定性，直接偏好优化（Direct Preference Optimization, DPO）等方法在不显式构建奖励模型的前提下，直接利用偏好数据对模型参数进行优化。该类方法通过将偏好关系转化为对数似然约束，实现对生成策略的直接调整，从而在一定程度上简化了对齐流程并提高了训练效率。尽管如此，其本质仍依赖于有限的偏好样本分布，因此同样难以完全覆盖开放环境中的复杂输入空间\cite{rafailov2023dpo,hong2024orpo,ethayarajh2024kto}。

除上述内部对齐方法外，安全对齐体系还依赖于模型外部的安全增强机制。在实际部署过程中，常通过输入检测、提示重写及输出审查等方式对模型行为进行二次约束，以降低越狱攻击带来的风险。此外，红队测试（red teaming）作为重要的评估与增强手段，通过构造对抗性输入系统性探测模型的潜在漏洞，为对齐策略的迭代优化提供依据\cite{rebedea2023nemo,markov2022holistic,yu2024gptfuzzer}。

总体而言，当前大语言模型的安全对齐技术体系呈现出“内部训练约束—外部防护增强—持续评估优化”的协同结构。SFT、RLHF及DPO等方法在常规场景下能够显著提升模型的安全表现，但其约束机制主要建立在有限样本与偏好近似之上，对分布外输入及对抗性攻击的鲁棒性仍有待进一步提升。这种以内在统计学习为基础的对齐方式，也为后续越狱攻击提供了可利用空间。

\begin{figure}[htbp]
\centering
\resizebox{0.98\textwidth}{!}{\begin{tikzpicture}[
    x=1pt,
    y=1pt,
    font=\small,
    >=stealth,
    panel/.style={draw=gray!40, rounded corners=14pt, fill=gray!4, line width=1pt},
    block/.style={draw=#1!75!black, fill=#1!11, rounded corners=9pt, line width=0.9pt, align=center},
    tag/.style={draw=#1!70!black, fill=white, rounded corners=6pt, line width=0.8pt, inner xsep=6pt, inner ysep=3pt, font=\scriptsize},
    flow/.style={->, line width=0.95pt, draw=gray!75},
    auxflow/.style={->, line width=0.85pt, draw=teal!70!black, dashed}
]
\def\W{730}
\def\H{268}
\fill[gray!4, rounded corners=16pt] (0,0) rectangle (\W,\H);
\draw[gray!35, rounded corners=16pt, line width=1pt] (0,0) rectangle (\W,\H);
\fill[cyan!18, rounded corners=14pt] (18,220) rectangle (712,250);
\node[anchor=west, font=\bfseries\small] at (34,235) {Safety Alignment Architecture};
\node[anchor=east, font=\scriptsize, text=gray!70!black] at (694,235) {internal alignment + external safeguards + continuous evaluation};

\draw[panel] (46,36) rectangle (548,192);
\node[anchor=west, font=\bfseries\scriptsize, text=blue!60!black] at (62,176) {内部对齐路径};

\node[block=gray, minimum width=132pt, minimum height=40pt] (pretrain) at (296,145) {\textbf{预训练模型}\\基础语言能力};
\node[block=blue, minimum width=110pt, minimum height=50pt] (sft) at (132,92) {\textbf{SFT}\\监督微调\\安全响应模式};
\node[block=teal, minimum width=110pt, minimum height=50pt] (rlhf) at (296,92) {\textbf{RLHF}\\奖励建模\\偏好优化};
\node[block=orange, minimum width=110pt, minimum height=50pt] (dpo) at (460,92) {\textbf{DPO}\\直接偏好优化\\分布调整};
\node[tag=blue] at (132,53) {训练阶段};
\node[tag=teal] at (296,53) {偏好校正};
\node[tag=orange] at (460,53) {稳定优化};

\draw[flow] (pretrain.south west) -- (sft.north);
\draw[flow] (pretrain.south) -- (rlhf.north);
\draw[flow] (pretrain.south east) -- (dpo.north);

\draw[panel] (566,36) rectangle (684,192);
\node[anchor=west, font=\bfseries\scriptsize, text=red!60!black] at (578,176) {外部防护机制};
\node[block=red, minimum width=86pt, minimum height=28pt] (inputf) at (625,145) {\textbf{输入过滤}};
\node[block=violet, minimum width=86pt, minimum height=28pt] (rewrite) at (625,111) {\textbf{提示重写}};
\node[block=magenta, minimum width=86pt, minimum height=28pt] (audit) at (625,77) {\textbf{输出审查}};
\node[block=brown, minimum width=86pt, minimum height=28pt] (redteam) at (625,43) {\textbf{红队测试}};

\node[block=cyan, minimum width=168pt, minimum height=32pt] (deploy) at (296,18) {\textbf{部署阶段模型服务}};
\draw[flow] (sft.south) -- (deploy.north west);
\draw[flow] (rlhf.south) -- (deploy.north);
\draw[flow] (dpo.south) -- (deploy.north east);

\draw[auxflow] (deploy.east) .. controls (420,18) and (560,145) .. (inputf.west);
\draw[auxflow] (deploy.east) .. controls (430,24) and (560,111) .. (rewrite.west);
\draw[auxflow] (deploy.east) .. controls (430,30) and (560,77) .. (audit.west);
\draw[auxflow] (redteam.west) .. controls (540,43) and (520,164) .. node[pos=0.45, above, font=\scriptsize, text=gray!70!black] {评估反馈} (pretrain.east);

\node[font=\scriptsize, text=gray!75!black, align=left] at (170,18) {协同关系: 内部训练塑造行为边界, 外部机制在部署时进行二次约束};
\end{tikzpicture}
}
\caption{安全对齐技术体系总览图}
\label{fig:safety-alignment-overview}
\end{figure}

\begin{table}[htbp]
\centering
\footnotesize
\setlength{\tabcolsep}{4pt}
\renewcommand{\arraystretch}{0.95}
\caption{第二章表格占位表：主流安全对齐方法比较表}
\begin{tabular}{p{0.15\textwidth}p{0.24\textwidth}p{0.24\textwidth}p{0.25\textwidth}}
\toprule
\textbf{方法} & \textbf{核心目标} & \textbf{优势} & \textbf{局限} \\
\midrule
SFT & 拟填“基于监督样本学习安全响应模式” & 拟填“训练稳定、实现简单、行为边界清晰” & 拟填“依赖样本覆盖，泛化有限” \\
RLHF & 拟填“基于奖励模型优化有用性与安全性平衡” & 拟填“能进行全局偏好调节” & 拟填“奖励偏移、训练复杂、对分布外样本敏感” \\
DPO & 拟填“基于偏好对直接调整输出分布” & 拟填“流程较简洁、优化较稳定” & 拟填“仍依赖偏好样本分布，缺乏硬约束” \\
\bottomrule
\end{tabular}
\end{table}
\noindent\textbf{附注：}该表建议直接作为论文正式表格，不需要额外绘图；内容可结合参考文献中的方法介绍和你第二章的文字总结进一步压缩润色。

\section{主流对齐方法及其作用机理}

当前大语言模型的安全对齐主要通过后训练阶段实现，其中以有监督微调（SFT）、基于人类反馈的强化学习（RLHF）以及直接偏好优化（DPO）为代表的方法构成了主流技术路径\cite{rafailov2023dpo,hong2024orpo,ethayarajh2024kto}。这些方法在形式上有所差异，但其本质均是通过引入外部偏好信息，对模型的生成分布进行重塑，从而改变其在特定输入条件下的输出行为。

从机理上看，大语言模型的生成过程可以视为对条件概率分布$P_\theta(y|x)$的建模，而安全对齐的核心目标在于对该分布进行约束，使其在满足任务需求的同时，降低生成不安全内容的概率。不同对齐方法的差异主要体现在约束施加的方式及优化路径上。

有监督微调（SFT）通过构建高质量的指令—响应数据，对预训练模型进行进一步训练，其优化目标通常仍为最大似然估计。在该过程中，模型通过模仿数据中的“安全响应模式”，逐步调整其条件分布，使得在相似输入下更倾向于生成符合规范的输出。从机制上看，SFT本质上是一种基于行为模仿的分布调整方法，其约束作用来源于示例数据对模型输出空间的重塑。然而，由于其优化目标仍然是局部条件概率的最大化，模型并未显式学习安全约束本身，而是学习其在训练数据中的表征形式。

在SFT的基础上，基于人类反馈的强化学习（RLHF）通过引入奖励模型，对生成结果进行全局质量评估。具体而言，RLHF首先利用人类偏好数据训练奖励模型$R_\phi(x, y)$，用于衡量输出的优劣；随后通过强化学习算法（如PPO）优化语言模型参数，使其生成结果在奖励函数下具有更高期望值。从优化角度看，RLHF将原始的似然最大化问题转化为期望奖励最大化问题，使模型在候选输出之间进行策略选择，从而在一定程度上实现“有用性”与“安全性”的权衡。然而，奖励模型本质上仍是对有限偏好数据的近似，其评估能力受数据分布与标注一致性的限制，可能在复杂或未见场景中产生偏差。

直接偏好优化（DPO）则在不显式构建奖励模型的情况下，直接利用成对偏好数据对模型参数进行优化。其核心思想是将人类偏好关系转化为对数似然约束，使模型在给定输入下对“优选响应”的概率高于“劣选响应”。与RLHF相比，DPO避免了奖励模型训练及强化学习过程中的不稳定性问题，在优化路径上更加简洁。从机理上看，DPO仍然通过调整条件概率分布来体现偏好信息，其约束作用直接来源于偏好数据本身，而非中间奖励函数。

综合来看，SFT、RLHF与DPO在形式上分别对应行为模仿、奖励驱动与偏好约束三种不同的对齐路径，但其共同特征在于：均依赖有限数据对模型分布进行统计性修正，而未引入显式的安全规则或可验证约束。这种基于数据驱动的对齐机制在常见场景下能够有效引导模型行为，但在面对分布外输入或对抗性构造时，其约束能力可能出现不稳定性。

此外，这些方法在训练过程中通常假设输入分布与实际应用环境一致，而在开放场景下，用户输入具有高度不确定性与可操控性。攻击者可以通过对输入进行语义重构、上下文扩展或结构扰动，使其偏离训练分布，从而削弱对齐机制的约束效果。这一现象表明，现有对齐方法在本质上仍属于基于经验分布的软约束，其安全边界缺乏稳定的理论保证。

\begin{equation}
\max_{\pi_\theta} \mathbb{E}_{y \sim \pi_\theta(\cdot \mid x)}\left[R_\phi(x,y)\right]-\beta \, \mathrm{KL}\!\left(\pi_\theta \,\|\, \pi_{\mathrm{ref}}\right)
\end{equation}
其中，$\pi_\theta$表示待优化策略，$R_\phi(x,y)$表示奖励模型对输出质量的评估结果，$\pi_{\mathrm{ref}}$表示参考模型，$\beta$表示KL正则项的权重系数。该式刻画了RLHF在奖励最大化与策略偏移控制之间的优化平衡。

\begin{equation}
\log \sigma \left(\beta \log \frac{\pi_\theta(y_w \mid x)}{\pi_{\mathrm{ref}}(y_w \mid x)}-\beta \log \frac{\pi_\theta(y_l \mid x)}{\pi_{\mathrm{ref}}(y_l \mid x)}\right)
\end{equation}
其中，$y_w$表示优选响应，$y_l$表示劣选响应，$\sigma(\cdot)$表示Sigmoid函数。该式说明DPO通过直接拉大优选响应与劣选响应之间的相对对数概率差，实现对偏好关系的显式建模。

\section{安全对齐中的目标冲突问题}

大语言模型的安全对齐过程本质上可视为一个多目标优化问题。在实际应用中，模型不仅需要具备较强的任务完成能力（helpfulness），还需满足安全性约束（harmlessness）及合规性要求。这些目标在多数情况下可以协同优化，但在复杂或边界情境下往往存在潜在冲突，从而导致模型行为的不稳定性。

从优化角度来看，预训练阶段的语言模型主要通过最大化条件概率分布来提升生成能力，其目标是提高输出的相关性与信息丰富性。而在对齐阶段，引入的人类偏好与安全约束则倾向于抑制某些高概率但潜在不安全的生成路径。因此，对齐过程实际上是在原始生成目标与安全约束之间进行权衡，即在“生成尽可能有用的内容”与“避免生成潜在风险内容”之间寻找平衡点。

这种目标之间的张力在具体任务中表现为不一致的优化方向。例如，在某些敏感问题场景中，提供详细且有用的信息可能同时增加潜在风险，而严格的安全约束则可能导致模型过度拒答，降低其实用性。为缓解这一问题，RLHF等方法通常通过奖励函数对不同目标进行加权整合，使模型在期望意义下取得折中解。然而，该折中结果依赖于奖励设计及偏好数据分布，在不同输入条件下可能呈现出不稳定行为\cite{rafailov2023dpo,ethayarajh2024kto}。

进一步地，由于对齐过程缺乏显式的约束边界，其优化结果通常以隐式方式编码于模型参数之中。这种基于统计学习的约束方式难以在所有输入空间中保持一致性，尤其是在分布外或对抗性输入条件下，模型可能优先满足“有用性”目标，从而弱化安全约束，产生对齐失效现象。

此外，不同对齐方法在处理目标冲突时采用的策略存在差异。SFT主要通过示例数据隐式传递行为偏好，其对冲突情境的处理能力有限；RLHF通过奖励模型对输出进行全局评估，在一定程度上缓解了目标冲突，但其效果依赖于奖励模型的准确性与稳定性；DPO则通过偏好对约束生成分布，其优化结果同样受限于偏好数据的表达能力。这些方法均未从机制上消除目标冲突，而是通过数据驱动的方式进行近似协调。

综上，安全对齐中的目标冲突问题是导致模型行为不稳定的重要因素之一。在缺乏统一优化目标及明确约束边界的情况下，模型在不同输入条件下可能表现出差异化的行为倾向。这种内在冲突为越狱攻击提供了可利用空间，即通过构造特定输入，使模型在“有用性”与“安全性”的权衡过程中偏向前者，从而绕过既有安全约束。

\begin{equation}
\max_{\theta} \ \lambda_1 U(\theta; x)-\lambda_2 H(\theta; x)
\end{equation}
其中，$U(\theta; x)$表示模型在输入$x$上的有用性目标，$H(\theta; x)$表示潜在风险或有害性目标，$\lambda_1$与$\lambda_2$分别表示两类目标的权重系数。该式用于说明安全对齐实质上是在有用性与安全性之间求取折中解的多目标优化过程。

\begin{figure}[htbp]
\centering
\resizebox{0.98\textwidth}{!}{\begin{tikzpicture}[
    x=1pt,
    y=1pt,
    font=\small,
    >=stealth,
    axis/.style={->, line width=1pt, draw=gray!80},
    grid/.style={draw=gray!20, line width=0.5pt},
    point/.style={circle, draw=#1!70!black, fill=#1!22, line width=1pt, minimum size=14pt, inner sep=0pt},
    bubble/.style={draw=#1!70!black, fill=#1!10, rounded corners=8pt, line width=0.85pt, align=left, inner xsep=8pt, inner ysep=5pt},
    helper/.style={draw=gray!55, dashed, line width=0.8pt}
]
\def\W{730}
\def\H{252}
\fill[gray!4, rounded corners=16pt] (0,0) rectangle (\W,\H);
\draw[gray!35, rounded corners=16pt, line width=1pt] (0,0) rectangle (\W,\H);
\fill[cyan!18, rounded corners=14pt] (18,204) rectangle (712,234);
\node[anchor=west, font=\bfseries\small] at (34,219) {Safety-Usefulness Trade-off};
\node[anchor=east, font=\scriptsize, text=gray!70!black] at (694,219) {multi-objective tension under alignment};

\foreach \x in {108,180,252,324,396,468,540,612}
  \draw[grid] (\x,48) -- (\x,176);
\foreach \y in {64,96,128,160}
  \draw[grid] (92,\y) -- (628,\y);

\draw[axis] (92,48) -- (640,48);
\draw[axis] (92,48) -- (92,184);
\node[anchor=north, font=\scriptsize] at (366,34) {任务完成度 / 有用性};
\node[anchor=south, rotate=90, font=\scriptsize] at (58,118) {安全遵循度};

\fill[green!7, rounded corners=12pt] (468,126) rectangle (620,174);
\node[font=\scriptsize, text=green!45!black] at (544,166) {理想平衡区};

\node[point=green] (benign) at (520,148) {};
\node[bubble=green, minimum width=112pt] at (575,144) {\textbf{正常良性请求}\\高完成度\\高安全遵循};

\node[point=orange] (wrapped) at (394,104) {};
\node[bubble=orange, minimum width=122pt] at (478,96) {\textbf{语义包装请求}\\看似合理\\边界开始松动};

\node[point=red] (danger) at (548,68) {};
\node[bubble=red, minimum width=120pt] at (586,72) {\textbf{直接危险请求}\\高风险\\应触发拒答};

\draw[helper] (wrapped) -- (wrapped |- 92,48);
\draw[helper] (wrapped) -- (wrapped -| 92,48);
\draw[helper] (danger) -- (danger |- 92,48);
\draw[helper] (benign) -- (benign -| 92,48);

\draw[->, line width=1pt, draw=violet!70!black] (420,118) .. controls (470,122) and (500,133) .. node[pos=0.55, above, font=\scriptsize] {对齐期望方向} (520,148);
\draw[->, line width=1pt, draw=red!70!black] (432,102) .. controls (466,84) and (520,72) .. node[pos=0.55, below, font=\scriptsize] {攻击诱导方向} (548,68);

\node[bubble=blue, minimum width=164pt] at (180,144) {\textbf{解释}\\对齐需要在“尽量回答”与“保持安全”之间取得折中,\\不同输入会把模型推向不同区域};
\end{tikzpicture}
}
\caption{安全性与有用性目标冲突示意图}
\label{fig:objective-conflict}
\end{figure}

\section{对齐表层性与行为泛化局限}

现有大语言模型的安全对齐方法主要依赖数据驱动的统计学习机制，其约束作用通常体现在对特定输入—输出模式的学习上，而非对安全语义的深层理解。这一特性使得对齐结果在很大程度上表现为“表层对齐”（superficial alignment），即模型倾向于学习可观测的行为模式，而非内在的安全原则。

从机制上看，无论是SFT、RLHF还是DPO，其优化过程均以训练数据中的示例或偏好为依据，对模型条件概率分布进行调整。在这一过程中，模型通过捕捉输入与输出之间的统计相关性，学习在特定语境下应当采取的响应策略。例如，在涉及敏感请求时，模型可能学习到以固定形式进行拒答或提供安全提示。然而，这种学习方式并未使模型形成对“何为不安全”的抽象语义表示，而是将安全约束映射为一类特定的语言表达模式。

由于缺乏语义层面的约束能力，当输入形式发生变化时，模型已学习的行为模式可能无法正确触发。例如，通过对问题进行语义重写、编码转换或上下文包装，攻击者可以构造与训练分布差异较大的输入，使其规避模型已习得的拒答模式\cite{jiang2024artprompt,deng2024multilingual,zeng2024johnny}。在此情况下，模型仍可能沿着原有的语言建模路径生成高概率但不安全的内容，从而表现出对齐失效。

进一步地，这种表层对齐特性还会受到分布外输入的影响。在开放环境中，用户输入具有高度多样性，而训练数据只能覆盖有限的表达形式。当输入超出训练分布时，模型往往依赖其通用语言能力进行推断，而非遵循对齐阶段学习到的行为模式。这种“泛化偏移”使得模型在复杂语境或组合型输入条件下更容易出现不符合预期的输出行为。

此外，在多轮对话与长上下文场景中，模型的行为不仅取决于单轮输入，还受到历史上下文的累积影响。当对齐约束未能在整个上下文空间中保持一致时，模型可能在局部语境中偏离既定的安全策略。这种现象进一步放大了表层对齐带来的局限，使得攻击者能够通过逐步引导的方式削弱模型的安全约束\cite{anil2024manyshot,chao2024pair}。

\begin{figure}[htbp]
\centering
\resizebox{0.98\textwidth}{!}{\begin{tikzpicture}[
    x=1pt,
    y=1pt,
    font=\small,
    >=stealth,
    panel/.style={draw=gray!40, rounded corners=14pt, fill=gray!4, line width=1pt},
    card/.style={draw=#1!72!black, fill=#1!10, rounded corners=9pt, line width=0.9pt, align=center},
    flow/.style={->, line width=0.95pt, draw=gray!75},
    weak/.style={->, line width=0.85pt, draw=red!70!black, dashed}
]
\def\W{730}
\def\H{262}
\fill[gray!4, rounded corners=16pt] (0,0) rectangle (\W,\H);
\draw[gray!35, rounded corners=16pt, line width=1pt] (0,0) rectangle (\W,\H);
\fill[cyan!18, rounded corners=14pt] (18,214) rectangle (712,244);
\node[anchor=west, font=\bfseries\small] at (34,229) {Superficial Alignment and Generalization Failure};
\node[anchor=east, font=\scriptsize, text=gray!70!black] at (694,229) {in-distribution refusal vs out-of-distribution drift};

\draw[panel] (36,34) rectangle (350,194);
\draw[panel] (380,34) rectangle (694,194);
\node[font=\bfseries\scriptsize, text=green!45!black] at (193,178) {训练分布内样本};
\node[font=\bfseries\scriptsize, text=red!55!black] at (537,178) {分布外变形样本};

\node[card=green, minimum width=94pt, minimum height=38pt] (in1) at (96,132) {显式危险提示};
\node[card=green, minimum width=94pt, minimum height=38pt] (in2) at (193,132) {拒答模式\\已学会触发};
\node[card=green, minimum width=94pt, minimum height=38pt] (in3) at (290,132) {安全输出\\边界清晰};
\draw[flow] (in1) -- (in2);
\draw[flow] (in2) -- (in3);
\node[card=green, minimum width=220pt, minimum height=34pt] at (193,76) {训练分布中的表层模式: 敏感词、固定拒答模板、显式风险信号};

\node[card=orange, minimum width=76pt, minimum height=34pt] (out1) at (438,132) {语义重写};
\node[card=orange, minimum width=76pt, minimum height=34pt] (out2) at (522,132) {编码转换};
\node[card=orange, minimum width=76pt, minimum height=34pt] (out3) at (606,132) {上下文包装};
\draw[flow] (out1) -- (out2);
\draw[flow] (out2) -- (out3);

\node[card=red, minimum width=214pt, minimum height=42pt] (blur) at (537,76) {拒答触发减弱\\模型转而依赖通用语言建模路径\\安全边界开始模糊};
\draw[weak] (out1.south) .. controls (438,102) and (490,88) .. (blur.west);
\draw[weak] (out2.south) -- (blur.north);
\draw[weak] (out3.south) .. controls (606,102) and (586,88) .. (blur.east);

\node[font=\scriptsize, text=gray!75!black, align=left] at (365,18) {结论: 对齐主要学习可观测行为模式, 当输入形式偏离训练分布时, 原有拒答模式可能不再稳定触发};
\end{tikzpicture}
}
\caption{表层对齐与泛化失效示意图}
\label{fig:superficial-alignment}
\end{figure}

综上，现有安全对齐方法在本质上是一种基于数据分布的表层行为约束，其有效性依赖于训练样本对输入空间的覆盖程度。在缺乏语义级约束与全局一致性机制的情况下，模型在面对分布外或对抗性输入时容易出现泛化失效。这种表层性与泛化局限构成了越狱攻击得以实现的重要基础。

\section{上下文与交互过程中的对齐失效}

大语言模型的生成行为具有显著的上下文依赖特性，其输出不仅取决于当前输入，还受到历史上下文的持续影响。在实际应用中，模型通常以多轮对话或长上下文形式运行，这使得安全对齐不再是针对单一输入的约束问题，而是需要在整个交互过程中保持一致性的动态约束。然而，现有对齐方法主要基于静态输入—输出样本进行训练，缺乏对复杂交互过程的全局建模能力，从而在多轮与长上下文场景下容易出现对齐失效。

从机制上看，大语言模型在生成过程中依赖条件概率分布$P(y|x_{1:n})$对当前上下文进行建模。当上下文中逐步引入新的信息时，模型的生成分布会随之发生动态变化。在这一过程中，早期对齐阶段所学习的安全行为模式可能被后续上下文所弱化或覆盖，导致模型输出逐渐偏离原有的安全约束。这种现象在长上下文条件下尤为明显，即随着上下文长度增加，模型对初始安全信号的依赖逐渐减弱。

在多轮交互场景中，对齐失效通常表现为累积式偏移。攻击者可以通过一系列看似合理的中间步骤逐步引导模型，使其在局部语境中接受某些隐含前提，从而在后续生成中输出不符合安全规范的内容。这种逐步诱导方式避免了直接触发显式的安全拒答机制，使攻击过程更加隐蔽且具有较高成功率\cite{chao2024pair,mehrotra2024tap,anil2024manyshot}。

此外，上下文结构的重构也是导致对齐失效的重要因素。例如，通过角色设定、情境假设或任务重定义等方式，可以改变模型对当前对话语境的理解，使其将原本受限的内容重新解释为“合理请求”\cite{zeng2024johnny,qi2024quack}。在此情况下，模型并非完全忽略安全约束，而是在新的语境下重新评估生成策略，从而产生与原始安全目标不一致的输出。

进一步地，在复杂系统场景中（如工具调用或智能体框架），模型的输入不再仅来自用户，还可能包含外部环境反馈或系统生成内容。这种多源上下文的叠加进一步扩大了输入空间的复杂性，使得对齐机制难以在所有路径上保持一致约束。在缺乏统一上下文建模与状态跟踪机制的情况下，模型行为更容易受到局部信息的干扰。

\begin{equation}
y_t \sim P_\theta(\cdot \mid x, h_{1:t-1})
\end{equation}
其中，$x$表示当前输入，$h_{1:t-1}$表示截至第$t-1$轮的历史上下文，$y_t$表示第$t$轮输出。该式表明多轮交互中的当前响应是由当前输入与历史对话共同决定的条件生成结果。

\begin{equation}
R_t=\alpha R_{t-1} + s(x_t,h_{1:t-1})
\end{equation}
其中，$R_t$表示第$t$轮的累计风险，$\alpha$表示历史风险的衰减系数，$s(x_t,h_{1:t-1})$表示当前轮输入在既有上下文条件下引入的风险增量。该式用于刻画多轮交互过程中风险随上下文逐步累积的动态特征。

\begin{figure}[htbp]
\centering
\resizebox{0.98\textwidth}{!}{\begin{tikzpicture}[
    x=1pt,
    y=1pt,
    font=\small,
    >=stealth,
    stage/.style={draw=#1!72!black, fill=#1!10, rounded corners=10pt, line width=0.95pt, minimum width=110pt, minimum height=54pt, align=center},
    risk/.style={draw=#1!72!black, fill=white, rounded corners=6pt, line width=0.8pt, inner xsep=6pt, inner ysep=3pt, font=\scriptsize},
    flow/.style={->, line width=1pt, draw=gray!78},
    trend/.style={line width=1.1pt, draw=red!70!black}
]
\def\W{730}
\def\H{248}
\fill[gray!4, rounded corners=16pt] (0,0) rectangle (\W,\H);
\draw[gray!35, rounded corners=16pt, line width=1pt] (0,0) rectangle (\W,\H);
\fill[cyan!18, rounded corners=14pt] (18,200) rectangle (712,230);
\node[anchor=west, font=\bfseries\small] at (34,215) {Multi-turn Alignment Failure Path};
\node[anchor=east, font=\scriptsize, text=gray!70!black] at (694,215) {risk accumulation across interaction rounds};

\node[stage=blue] (s1) at (92,112) {\textbf{初始试探}\\低风险提问\\边界摸索};
\node[stage=teal] (s2) at (226,112) {\textbf{角色巩固}\\设定身份\\重塑语境};
\node[stage=orange] (s3) at (360,112) {\textbf{规则豁免}\\弱化限制\\制造例外};
\node[stage=red!80!orange] (s4) at (494,112) {\textbf{任务细化}\\逐步补充\\逼近敏感目标};
\node[stage=red] (s5) at (628,112) {\textbf{最终突破}\\输出偏离\\安全失效};

\draw[flow] (s1) -- (s2);
\draw[flow] (s2) -- (s3);
\draw[flow] (s3) -- (s4);
\draw[flow] (s4) -- (s5);

\node[risk=blue] at (92,62) {风险低};
\node[risk=teal] at (226,62) {风险上升};
\node[risk=orange] at (360,62) {约束被松动};
\node[risk=red!80!orange] at (494,62) {累计偏移};
\node[risk=red] at (628,62) {高风险输出};

\draw[trend] (92,28) -- (226,40) -- (360,56) -- (494,74) -- (628,96);
\foreach \x/\y in {92/28,226/40,360/56,494/74,628/96}
  \fill[red!70!black] (\x,\y) circle (2.2pt);
\node[anchor=west, font=\scriptsize, text=red!70!black] at (78,18) {累计风险};

\node[font=\scriptsize, text=gray!72!black, align=left] at (360,18) {多轮交互的关键问题不在单次越界, 而在于历史上下文会持续重塑模型的条件分布};
\end{tikzpicture}
}
\caption{多轮交互中的对齐失效路径图}
\label{fig:multiturn-alignment-failure}
\end{figure}

综上，上下文依赖与交互动态性使得大语言模型的安全对齐问题从“静态约束”转变为“过程约束”。现有方法由于主要基于单轮样本进行训练，难以在多轮与复杂上下文环境中维持稳定的安全行为。这种在交互过程中的对齐失效，与目标冲突及表层对齐问题相互叠加，共同构成了越狱攻击得以实现的重要机制基础。

\section{本章小结}

本章围绕大语言模型的安全对齐机制及其潜在漏洞，从模型生成原理与对齐方法出发，对其内在机理进行了系统分析。首先，对大语言模型的基本建模方式进行了概述，指出其基于条件概率分布的生成机制决定了模型行为对数据分布与上下文输入的高度依赖。在此基础上，进一步梳理了当前主流的安全对齐技术体系，包括基于SFT、RLHF及DPO的内部训练方法，以及输入输出约束与红队测试等外部增强机制。

在对齐方法机理分析的基础上，本文从优化目标、行为表征及交互过程三个层面对安全对齐的局限性进行了深入探讨。具体而言，目标冲突问题揭示了模型在“有用性”与“安全性”之间难以实现统一优化；表层对齐特性表明模型主要学习的是可观测的行为模式，而非安全语义本身；上下文依赖与多轮交互则进一步放大了对齐约束的不稳定性，使模型在动态输入条件下更易出现偏离。

综合上述分析可以看出，现有安全对齐方法本质上仍属于基于数据分布的软约束，其有效性依赖于训练样本与偏好数据的覆盖范围，在面对分布外输入及对抗性构造时缺乏稳定的约束能力。上述机制性缺陷共同构成了越狱攻击得以实现的基础条件。

因此，有必要在统一分析框架下，对不同攻击方式在上述机制约束下的表现进行系统评估。基于此，下一章将围绕越狱攻击方法展开，通过构建统一实验框架，对不同攻击策略在多模型条件下的有效性进行分析，并进一步揭示模型安全对齐的脆弱性特征。
